\documentclass[a4paper,11pt]{article}
\usepackage[utf8]{inputenc}
\usepackage[T1]{fontenc}
\usepackage[french]{babel}
\usepackage{amsmath}
\usepackage[top=2cm, bottom=2cm, left=2.2cm, right=2.2cm]{geometry}
\usepackage{xcolor}
\usepackage{booktabs}
\usepackage{tabularx}
\usepackage{enumitem}
\usepackage{listings}
\usepackage{titlesec}
\usepackage{hyperref}
\hypersetup{colorlinks=true, linkcolor=primary, urlcolor=primary}

\definecolor{primary}{HTML}{1A237E}
\definecolor{success}{HTML}{2E7D32}
\definecolor{warning}{HTML}{E65100}
\definecolor{codebg}{HTML}{F5F5F5}

\lstset{backgroundcolor=\color{codebg}, basicstyle=\ttfamily\footnotesize,
  breaklines=true, frame=single, showstringspaces=false, columns=fullflexible,
  keywordstyle=\color{primary}}

\titleformat{\section}{\large\bfseries\color{primary}}{\thesection.}{0.6em}{}
\titlespacing*{\section}{0pt}{12pt}{4pt}
\titleformat{\subsection}{\bfseries}{\thesubsection.}{0.5em}{}
\setlength{\parindent}{0pt}
\newcommand{\code}[1]{\texttt{#1}}
\newcommand{\etape}[1]{\vspace{4pt}\textbf{\color{primary}#1}\par}

\title{\textbf{Démonstration V2V sur le terrain --- la voiture roule, le PC affiche son cap et sa vitesse}}
\author{Douae Sebti}
\date{9 juillet 2026}

\begin{document}
\maketitle
\vspace{-8pt}\hrule\vspace{10pt}

% =====================================================================
\section{Le but de la démonstration}
Je veux montrer ma contribution en conditions réelles : la voiture (Jetson) \textbf{roule},
et moi je la suis à pied avec le PC. Sur l'écran du PC s'affichent, \textbf{en direct}, la
\textbf{position}, le \textbf{cap} (direction) et la \textbf{vitesse} de la voiture ---
toutes issues de son odométrie et transportées dans les messages \textbf{CAM} ETSI que la
voiture diffuse par WiFi. Une \textbf{alerte de proximité} se déclenche quand la voiture est
à moins de 5~m. C'est la preuve concrète que le CAM porte maintenant le mouvement réel du
véhicule, et pas seulement une position.

% =====================================================================
\section{Avant de commencer (vérifications)}
\begin{itemize}[leftmargin=1.6em]
  \item Le PC \textbf{et} la Jetson sont sur le \textbf{même réseau WiFi} (le lien CAM utilise
        le multicast, qui reste local au WiFi).
  \item Le \textbf{LiDAR} est bien branché sur la Jetson : \code{ls /dev/ttyUSB0} et
        \code{lsusb | grep 10c4} doivent répondre (l'odométrie \emph{rf2o} vient du LiDAR).
  \item \code{socktap} est compilé des deux côtés (fait le 9~juillet).
  \item Trouver l'adresse IP de la Jetson (elle change avec le DHCP). Depuis le PC :
\begin{lstlisting}[language=bash]
ip neigh | grep -i 14:75:5b     # MAC de la Jetson -> son IP du jour (ex. 10.37.11.29)
\end{lstlisting}
\end{itemize}

% =====================================================================
\section{Étape 1 --- lancer la voiture (sur la Jetson)}
Depuis le PC, j'ouvre une session SSH vers la Jetson, puis je lance le script qui démarre
toute la chaîne (LiDAR + odométrie + conversion GPS + émission des CAM) :

\begin{lstlisting}[language=bash]
ssh protova2@10.37.11.29        # adapter l'IP au besoin
./launch_v2v_odom.sh
\end{lstlisting}

Ce script (domaine ROS 125) démarre :
\begin{itemize}[leftmargin=1.6em]
  \item le \textbf{LiDAR} + \code{rf2o} $\rightarrow$ odométrie \code{/odom\_rf2o} ;
  \item \code{odom\_to\_gps.py} : odométrie $\rightarrow$ \code{/v2v/my\_gps}
        \code{[lat, lon, cap, vitesse]} ;
  \item \code{v2v\_node.py} : pousse cet état à \code{socktap} $\rightarrow$ le \textbf{CAM
        émis} porte la position, le \textbf{cap} et la \textbf{vitesse} réels.
\end{itemize}
J'attends le message indiquant que \code{v2v\_node} est lancé (environ 15~s, le temps que
l'odométrie démarre). \textbf{Astuce sécurité :} pas besoin de mettre les moteurs (ESC) ;
je peux tout simplement \textbf{pousser la voiture à la main}, l'odométrie suit quand même.

% =====================================================================
\section{Étape 2 --- lancer le moniteur (sur le PC)}
Dans un terminal du PC, je lance le tableau de bord terrain :

\begin{lstlisting}[language=bash]
python3 ~/v2v_car_monitor.py --iface wlp2s0 \
  --my-lat 48.76680 --my-lon 11.43200 --alert 5
\end{lstlisting}

\begin{itemize}[leftmargin=1.6em]
  \item \code{--my-lat/--my-lon} = ma position de référence (par défaut, l'origine de
        l'odométrie de la voiture $\rightarrow$ au départ la distance est proche de 0).
  \item \code{--alert 5} = seuil d'alerte de proximité, en mètres.
  \item Le moniteur \textbf{n'a pas besoin de ROS} : il lance lui-même \code{socktap},
        reçoit les CAM de la voiture et affiche une ligne rafraîchie en continu.
\end{itemize}

% =====================================================================
\section{Étape 3 --- faire rouler et observer}
Je fais rouler (ou je pousse) la voiture. Sur le PC, la ligne se met à jour environ une fois
par seconde :

\begin{lstlisting}
[voiture #2] 48.766900,11.432100 | cap  45.0 NE / | v  3.00 m/s ( 10.8 km/h) | dist  13.3 m |   25 CAM @ 1.0 Hz
\end{lstlisting}

Ce qu'il faut regarder --- \textbf{c'est là qu'est la contribution} :
\begin{itemize}[leftmargin=1.6em]
  \item quand je \textbf{tourne} la voiture, le \textbf{cap} change (N, NE, E\ldots avec la
        flèche) ;
  \item quand je l'\textbf{accélère}, la \textbf{vitesse} (m/s et km/h) augmente ;
  \item la \textbf{position} et la \textbf{distance} évoluent en temps réel ;
  \item quand la voiture passe à moins de 5~m, l'\textbf{alerte} \code{*** ALERTE <5m ***}
        s'affiche.
\end{itemize}
Avant, ces champs cap/vitesse restaient à « indisponible » : les voir vivre prouve
l'enrichissement du CAM.

% =====================================================================
\section{Arrêt et dépannage}
\begin{itemize}[leftmargin=1.6em]
  \item \textbf{Arrêt} : \code{Ctrl-C} sur le moniteur (PC) et sur \code{launch\_v2v\_odom.sh}
        (Jetson).
  \item \textbf{Le moniteur affiche « indispo » une fois sur deux} : c'est qu'un ancien
        \code{socktap} (mode statique) tourne encore sur la Jetson et émet aussi en station~2.
        Le tuer : \code{ps -eo pid,args | grep "[s]ocktap"} puis \code{kill <PID>}
        (\code{launch\_v2v\_odom.sh} le fait normalement au démarrage).
  \item \textbf{Rien ne s'affiche} : vérifier que les deux machines sont sur le même WiFi, et
        que \code{socktap} tourne bien sur la Jetson (\code{ps ... | grep socktap}).
  \item \textbf{Le LiDAR se débranche} quand on manipule la voiture (connecteur CP2102) :
        rebrancher et vérifier \code{ls /dev/ttyUSB0} avant de relancer.
  \item \textbf{Chemin} : lancer \code{\char`~/vanetza/build/bin/socktap} (avec le \code{\char`~}),
        sinon « No such file or directory ».
\end{itemize}

% =====================================================================
\section{Ce que la démonstration prouve}
Deux véhicules (la voiture et le PC) échangent des messages \textbf{CAM ETSI} conformes,
sur WiFi, sans matériel radio spécialisé. Grâce à mon travail, le CAM émis par la voiture
porte désormais son \textbf{cap} et sa \textbf{vitesse} réels, calculés depuis l'odométrie,
et le PC les affiche en direct avec une alerte de proximité. C'est la brique de base d'un
système coopératif (V2V) : chaque véhicule sait où sont les autres, dans quelle direction et
à quelle vitesse ils vont. \textbf{Prochaine étape} : rejouer cette démonstration sur le
réseau \textbf{5G} (en passant le lien en unicast, le cellulaire bloquant le multicast).

\vspace{8pt}\hrule
{\footnotesize\itshape Fichiers : \code{v2v\_car\_monitor.py} (PC), \code{launch\_v2v\_odom.sh}
(Jetson), \code{v2v\_node.py}, \code{odom\_to\_gps.py}. Voir aussi
\code{seance\_cam\_v2v\_2026-07-09.tex} et \code{v2v\_explique\_pas\_a\_pas.tex}.}

\end{document}
